\section{The final comparison and proof of the theorem}\label{sec:conclusion}

\begin{proof}[Proof of \Cref{thm:main}]
Suppose bigeodesics exist with positive probability.  Each bigeodesic
contains a lattice vertex.  Countability of the vertices and the charts
in \Cref{lem:charts}, together with the exclusion of their deterministic
exceptional labels, yields a chart in which a bigeodesic through some
fixed vertex has positive probability.  Translate that vertex to zero
and use the lattice coordinates of the chart.  By \Cref{prop:cone}, its
arms are contained in $Q^\pm(0,H)$ for some finite $H$ almost surely.
Taking the countable union over integer $H$ gives a deterministic
$H_0\geq1$ for which the event $\mathcal B(0)$ in \Cref{sec:tilts} has
probability $p>0$.  Removing the countably many deterministic partition
endpoints does not change this probability, by the fixed-label
backward-finiteness assertion in \Cref{thm:structure}.

Apply \Cref{prop:tilt-marks} for each positive integer $M$, with the
fixed spacing $l$ from \eqref{eq:fixed-grid-spacing}.  We next impose
two filters on its marks.  The first supplies the base comparator rays
required for \Cref{lem:budget-comparison}; the second supplies their
required separation within each bin.

For a horizontal site $k$, let $\mathcal R_K(k)$ be the event that, in
the unchanged base environment, every ray from $k$ with label in $I$
lies in $Q^+(k,K)$ and every ray from $k$ with label in $-I$ lies in
$Q^-(k,K)$.  By \Cref{prop:cone},
\[
 \mu(\mathcal R_K(k)^c)\longrightarrow0\qquad(K\longrightarrow\infty).
\]
Choose one deterministic $K\geq1$, independent of $M$, such that this
failure probability is at most $p/8$.  It is the same at all horizontal
sites by translation invariance.  Retain a mark when
\[
 \widetilde J_M(k)=J_M(k)\mathbf1_{\mathcal R_K(k)}=1.
\]
Unconditional subtraction, without any independence assertion, gives
\begin{equation}\label{eq:retained-mark-probability}
 q_M:=\PP(\widetilde J_M(k)=1)
 \geq\PP(J_M(k)=1)-\mu(\mathcal R_K(k)^c)
 \geq p/8.
\end{equation}
This filter is imposed after selecting the paths.  Its complement has
an existential ray witness, so $\mathcal R_K(k)$ is a universally
measurable environment event by \Cref{lem:measurable}.  It remains a
separate filter; it is not placed inside an existential path projection.
In particular no new alternation of path quantifiers is used to define
the selected marks.  The translated rule continues to give equal
marginals along the grid.

The near radius $R_0$, spacing $l$, and comparator offset $K$ are fixed
independently of $M$. For each $M$ we choose $R(M)$ and the comparison
separation below, refine the partition, and only then increase $M$.
Fix $M$ for the moment. Its radius $R(M)$ is a fixed finite number.
The selected arms have offset $H_0$, the base comparator rays have
offset $K$, the changed region has radius at most $D_{R(M)}$ from
\eqref{eq:defect-radius}, and the prescribed gain portion has radius at
most $R(M)$.  The comparison in \Cref{lem:budget-comparison} therefore
holds at a finite deterministic horizontal separation $S(M)$.  The same
$S(M)$ works for every bin of every partition $\Pi_j$: the cone bounds
control all the relevant comparator labels in the chart, and the
geometric bounds on the changed region and gain portion are independent
of the bin.  Increase $S(M)$ if necessary to be nonnegative.  Its possible
growth with $M$ has no effect on the already fixed grid spacing $l$.

For the finite partition $\Pi_j$, retain a site with
$\widetilde J_M(k)=1$ only if there is no other site $y$ with
$\widetilde J_M(y)=1$ and
\[
 0<|y-k|\leq S(M)
\]
whose selected label lies in the same bin as $\rho_k$.  Here $k,y$ lie
on the horizontal grid, so the displayed distance is their horizontal
separation.  Denote this final indicator by $J_{M,j}(k)$ and its common
marginal probability by $q_{M,j}$.  The test concerns at most
$2\lfloor S(M)/l\rfloor$ other sites and is measurable from the
previously selected labels and marks.  At a fixed $j$, two surviving
marks in the same bin have separation strictly greater than $S(M)$,
so \Cref{lem:budget-comparison} applies to consecutive surviving marks
of that bin with gain $b=M$.

For completeness, the disappearance of this last filter is a
pointwise finite-set argument.  On the common full-measure event of
\Cref{sec:unions}, fix an outcome with $\widetilde J_M(k)=1$.  By
\Cref{prop:distinct-marks}, every marked site $y\ne k$, and in
particular every one of the finitely many retained neighbors within
distance $S(M)$, has $\rho_y\ne\rho_k$.  The nested partitions
eventually put each such pair in different bins.  The maximum of their
finitely many separating indices is finite, so $k$ survives the last
filter for every subsequent partition.  If $\widetilde J_M(k)=0$, all
of its final indicators are zero.  We have therefore proved
\begin{equation}\label{eq:fixed-M-filter-limit}
 J_{M,j}(k)\longrightarrow\widetilde J_M(k)
 \quad\text{almost surely as }j\longrightarrow\infty,
 \qquad q_{M,j}\longrightarrow q_M,
\end{equation}
where the probability convergence follows by bounded convergence.
No rate of separation of random labels is required.

For a bin $(\alpha,\beta)_I$ of $\Pi_j$, write $q_{M,j}^{\alpha,\beta}$
for the probability that a given site survives and its selected label
lies in that bin.  \Cref{prop:label-budget}, with $b=M$, gives
\[
 q_{M,j}^{\alpha,\beta}
 \leq\frac lM\bigl(r(\beta)-r(\alpha)\bigr).
\]
Every selected label avoids all partition endpoints, so each surviving
mark belongs to exactly one bin.  Summing over the finitely many bins
and using their total width $W_I$ yields
\begin{equation}\label{eq:final-finite-partition-budget}
 q_{M,j}\leq\frac{lW_I}{M}.
\end{equation}
The label-budget bound uses only equal marginal marking probabilities
and the pathwise comparison.  In particular the dependencies introduced
by the shared raises and both filters are permitted.

Now send $j$ to infinity for this fixed $M$.  Equations
\eqref{eq:retained-mark-probability}, \eqref{eq:fixed-M-filter-limit},
and \eqref{eq:final-finite-partition-budget} imply
\[
 \frac p8\leq q_M\leq\frac{lW_I}{M}.
\]
All of $p,l,W_I$ were fixed before $M$ was chosen.  Taking any integer
$M>8lW_I/p$ is a contradiction.

The reduction to a chart involved countably many vertices and charts.
Within the construction the excluded deterministic labels, integer
radii, sites, and union metrics form countable collections; equality of
the Busemann pairs and distinctness of marked labels hold pathwise for
their random labels.  Thus the contradiction excludes the existence
event itself and proves almost-sure absence of all bigeodesics, with no
prescribed deterministic direction.
\end{proof}

The perturbation argument uses only the nonatomic marginal law, two
separated sets of positive mass, and the structural conclusions
established earlier under the minimum-of-four second-moment assumption.
In particular all densities in the change-of-law calculation are taken
relative to $G$ itself, and the total gain is accumulated over finitely
many raised edges.  A density with respect to Lebesgue measure,
unbounded support, and a finite mean for one edge are not required.
